\documentclass[11pt,a4paper]{article}

% ===== PACKAGES =====
\usepackage[utf8]{inputenc}
\usepackage[T1]{fontenc}
\usepackage[margin=1in]{geometry}
\usepackage{hyperref}
\usepackage{listings}
\usepackage{xcolor}
\usepackage{booktabs}
\usepackage{array}
\usepackage{longtable}
\usepackage{fancyhdr}
\usepackage{titlesec}
\usepackage{tcolorbox}
\usepackage{fontawesome5}
\usepackage{amssymb}
\usepackage{amsmath}
\usepackage{enumitem}

% ===== COLORS =====
\definecolor{codegreen}{rgb}{0.25,0.49,0.48}
\definecolor{codegray}{rgb}{0.5,0.5,0.5}
\definecolor{codepurple}{rgb}{0.58,0,0.82}
\definecolor{codeblue}{rgb}{0.13,0.29,0.53}
\definecolor{codeorange}{rgb}{0.8,0.4,0.0}
\definecolor{backcolour}{rgb}{0.95,0.95,0.95}
\definecolor{warningbg}{rgb}{1.0,0.95,0.85}
\definecolor{warningborder}{rgb}{0.9,0.6,0.2}
\definecolor{tipbg}{rgb}{0.9,0.97,0.9}
\definecolor{tipborder}{rgb}{0.3,0.7,0.3}
\definecolor{notebg}{rgb}{0.9,0.93,1.0}
\definecolor{noteborder}{rgb}{0.3,0.4,0.7}

% ===== CODE LISTING STYLE =====
\lstdefinelanguage{IEC61131}{
    keywords={VAR, END_VAR, VAR_INPUT, VAR_OUTPUT, VAR_IN_OUT, VAR_GLOBAL, 
              TYPE, END_TYPE, STRUCT, END_STRUCT, FUNCTION, END_FUNCTION,
              FUNCTION_BLOCK, END_FUNCTION_BLOCK, PROGRAM, END_PROGRAM,
              IF, THEN, ELSE, ELSIF, END_IF, CASE, OF, END_CASE,
              FOR, TO, BY, DO, END_FOR, WHILE, END_WHILE, REPEAT, UNTIL, END_REPEAT,
              TRUE, FALSE, AND, OR, NOT, XOR, MOD,
              ARRAY, OF, POINTER, REFERENCE, TO, AT, EXTENDS, IMPLEMENTS,
              METHOD, END_METHOD, PROPERTY, END_PROPERTY,
              PUBLIC, PRIVATE, PROTECTED, INTERNAL, FINAL, ABSTRACT,
              THIS, SUPER, VAR_STAT, CONSTANT},
    keywordstyle=\color{codeblue}\bfseries,
    keywords=[2]{BOOL, BYTE, WORD, DWORD, LWORD, SINT, USINT, INT, UINT, 
                 DINT, UDINT, LINT, ULINT, REAL, LREAL, STRING, WSTRING,
                 TIME, LTIME, DATE, TIME_OF_DAY, TOD, DATE_AND_TIME, DT,
                 TcEventSeverity, T_FILETIME},
    keywordstyle=[2]\color{codepurple}\bfseries,
    keywords=[3]{ADR, SIZEOF, REF, FILETIME_TO_DT, GETSYSTEMTIME},
    keywordstyle=[3]\color{codeorange}\bfseries,
    sensitive=false,
    morecomment=[l]{//},
    morecomment=[s]{(*}{*)},
    commentstyle=\color{codegreen}\itshape,
    stringstyle=\color{codeorange},
    morestring=[b]',
    morestring=[b]",
}

\lstdefinestyle{iecstyle}{
    language=IEC61131,
    backgroundcolor=\color{backcolour},
    basicstyle=\ttfamily\small,
    breakatwhitespace=false,
    breaklines=true,
    captionpos=b,
    keepspaces=true,
    numbers=left,
    numbersep=8pt,
    numberstyle=\tiny\color{codegray},
    showspaces=false,
    showstringspaces=false,
    showtabs=false,
    tabsize=4,
    frame=single,
    framerule=0.5pt,
    rulecolor=\color{codegray},
    xleftmargin=15pt,
    framexleftmargin=15pt,
}

\lstset{style=iecstyle}

% ===== TCOLORBOX STYLES =====
\tcbuselibrary{skins,breakable}

\newtcolorbox{warningbox}{
    colback=warningbg,
    colframe=warningborder,
    boxrule=1pt,
    left=10pt,
    right=10pt,
    top=8pt,
    bottom=8pt,
    breakable,
    before upper={\faExclamationTriangle\ \textbf{Warning:} }
}

\newtcolorbox{tipbox}{
    colback=tipbg,
    colframe=tipborder,
    boxrule=1pt,
    left=10pt,
    right=10pt,
    top=8pt,
    bottom=8pt,
    breakable,
    before upper={\faLightbulb\ \textbf{Tip:} }
}

\newtcolorbox{notebox}{
    colback=notebg,
    colframe=noteborder,
    boxrule=1pt,
    left=10pt,
    right=10pt,
    top=8pt,
    bottom=8pt,
    breakable,
    before upper={\faInfoCircle\ \textbf{Note:} }
}

% ===== HEADER/FOOTER =====
\pagestyle{fancy}
\fancyhf{}
\fancyhead[L]{\leftmark}
\fancyhead[R]{TwinCAT 3 Lecture Notes}
\fancyfoot[C]{\thepage}
\renewcommand{\headrulewidth}{0.4pt}
\renewcommand{\footrulewidth}{0.4pt}

% ===== HYPERREF SETUP =====
\hypersetup{
    colorlinks=true,
    linkcolor=codeblue,
    filecolor=magenta,
    urlcolor=cyan,
    pdftitle={TwinCAT 3 - Function Blocks and OOP},
    pdfauthor={},
    bookmarks=true,
}

% ===== TITLE FORMATTING =====
\titleformat{\section}
    {\normalfont\Large\bfseries}{\thesection}{1em}{}[{\titlerule[0.8pt]}]
\titleformat{\subsection}
    {\normalfont\large\bfseries}{\thesubsection}{1em}{}
\titleformat{\subsubsection}
    {\normalfont\normalsize\bfseries}{\thesubsubsection}{1em}{}

% ===== DOCUMENT INFO =====
\title{
    \vspace{-1cm}
    \textbf{Lecture Notes: TwinCAT 3}\\
    \Large Function Blocks and Object-Oriented Programming (OOP)
}
\author{}
\date{\today}

% ===== BEGIN DOCUMENT =====
\begin{document}

\maketitle

\tableofcontents
\newpage

% ============================================================
\section{Procedural vs. Object-Oriented Programming}
% ============================================================

The transition from functions to \textbf{function blocks (FBs)} represents a shift from \textbf{Procedural Programming (POP)} to \textbf{Object-Oriented Programming (OOP)}.

\begin{itemize}[leftmargin=*]
    \item \textbf{Procedural (POP):} Focuses on a sequence of actions. Global data is defined, and functions act upon that data.
    
    \item \textbf{Object-Oriented (OOP):} Focuses on data. Programs are divided into \textbf{objects} that carry their own state (data) and behavior (methods).
    
    \item \textbf{The Object:} In industrial automation, an object can represent physical hardware (a motor, a sensor) or abstract concepts (a database logger, a mathematical calculation).
\end{itemize}

\begin{table}[h]
\centering
\begin{tabular}{@{}lll@{}}
\toprule
\textbf{Aspect} & \textbf{Procedural (POP)} & \textbf{Object-Oriented (OOP)} \\
\midrule
Focus & Sequence of actions & Data and objects \\
Data Location & Global variables & Encapsulated in objects \\
Code Reuse & Functions & Classes/Function Blocks \\
State Management & External to functions & Internal to objects \\
\bottomrule
\end{tabular}
\caption{Comparison of Programming Paradigms}
\end{table}

% ============================================================
\section{Function Blocks as Blueprints}
% ============================================================

A function block serves as the \textbf{blueprint} for an object.

\begin{itemize}[leftmargin=*]
    \item \textbf{State:} Defined by internal variables (e.g., a motor's current velocity or temperature).
    
    \item \textbf{Behavior:} Defined by \textbf{methods} (e.g., \texttt{PowerEnable}, \texttt{SetSpeed}).
    
    \item \textbf{Instantiation:} To use an FB, you must create an \textbf{instance} of it. This reserves memory in the PLC, allowing multiple objects (e.g., \texttt{fbWinch} and \texttt{fbConveyor}) to exist independently, each holding its own unique data.
\end{itemize}

\subsection*{Example: Instantiation}

\begin{lstlisting}
// Program: MAIN
// Declaring multiple instances of the Motor function block
// Each instance maintains its own independent state and data
PROGRAM MAIN
VAR
    fbWinch    : FB_Motor;    // First instance - controls the winch motor
    fbConveyor : FB_Motor;    // Second instance - controls the conveyor motor
    fbPump     : FB_Motor;    // Third instance - controls the pump motor
END_VAR

// Each instance can be controlled independently
fbWinch(bEnable := TRUE, fTargetSpeed := 50.0);
fbConveyor(bEnable := TRUE, fTargetSpeed := 75.0);
fbPump(bEnable := FALSE, fTargetSpeed := 0.0);

// Access instance-specific data
// fbWinch.fCurrentSpeed will be different from fbConveyor.fCurrentSpeed
\end{lstlisting}

\begin{notebox}
Each instance of a function block has its own memory allocation. This means \texttt{fbWinch} and \texttt{fbConveyor} can have completely different internal states even though they are created from the same blueprint.
\end{notebox}

% ============================================================
\section{The Pillars of OOP in TwinCAT 3}
% ============================================================

\subsection{Encapsulation}

Bundling data with the methods that operate on it and hiding the internal state from the outside. This prevents ``spaghetti code'' caused by excessive global variables.

\subsection{Abstraction}

Hiding complex internal logic from the user. The caller only needs to know \textit{what} a method does (e.g., \texttt{AddEvent}), not the specific details of \textit{how} it is accomplished.

\subsection{Inheritance}

The process where a \textbf{child function block} inherits all attributes and methods of a \textbf{parent function block}.

\begin{itemize}[leftmargin=*]
    \item The child can add its own specific features or methods.
    \item \textbf{``IS-A'' Relationship:} An SMS notification \textit{is a} notification.
    \item \textbf{Syntax:} Use the \texttt{EXTENDS} keyword to implement inheritance.
\end{itemize}

\subsection*{Example: Inheritance}

\begin{lstlisting}
// Parent Function Block: FB_Notification
// Base class containing common notification functionality
FUNCTION_BLOCK FB_Notification
VAR
    sRecipient   : STRING(255);    // Who receives the notification
    sMessage     : STRING(1000);   // The notification content
    dtTimestamp  : DATE_AND_TIME;  // When the notification was created
    bSent        : BOOL;           // Whether notification has been sent
END_VAR

// Methods for the parent FB would be defined here
// e.g., Send(), GetStatus(), etc.
\end{lstlisting}

\begin{lstlisting}
// Child Function Block: FB_SMSNotification
// Inherits from FB_Notification and adds SMS-specific functionality
FUNCTION_BLOCK FB_SMSNotification EXTENDS FB_Notification
VAR
    sPhoneNumber    : STRING(20);     // SMS-specific: recipient phone number
    nMaxCharacters  : INT := 160;     // SMS-specific: character limit
    bDeliveryReport : BOOL;           // SMS-specific: request delivery confirmation
END_VAR

// This FB automatically has access to:
// - sRecipient (inherited)
// - sMessage (inherited)
// - dtTimestamp (inherited)
// - bSent (inherited)
// Plus its own SMS-specific variables
\end{lstlisting}

\begin{lstlisting}
// Another Child: FB_EmailNotification
// Different specialization of the base notification class
FUNCTION_BLOCK FB_EmailNotification EXTENDS FB_Notification
VAR
    sEmailAddress : STRING(255);      // Email-specific: recipient address
    sSubject      : STRING(255);      // Email-specific: email subject line
    sAttachment   : STRING(255);      // Email-specific: attachment path
    bHTMLFormat   : BOOL;             // Email-specific: HTML or plain text
END_VAR
\end{lstlisting}

\begin{tipbox}
Use inheritance when you have a clear ``IS-A'' relationship. For example, \texttt{FB\_SMSNotification} IS-A \texttt{FB\_Notification}. This promotes code reuse and maintains a logical hierarchy.
\end{tipbox}

% ============================================================
\section{Methods and Access Specifiers}
% ============================================================

Methods are used to change or access the internal state of an object. Every method requires an \textbf{access specifier} to define who can call it:

\begin{table}[h]
\centering
\begin{tabular}{@{}llp{7cm}@{}}
\toprule
\textbf{Specifier} & \textbf{Symbol} & \textbf{Accessibility} \\
\midrule
\texttt{PUBLIC} & + & Accessible by everyone (the FB itself and external code) \\
\texttt{PRIVATE} & - & Accessible only within the FB where it is declared \\
\texttt{PROTECTED} & \# & Accessible by the FB and its children (via inheritance) \\
\texttt{INTERNAL} & \texttildelow & Accessible within the FB and its local library \\
\bottomrule
\end{tabular}
\caption{Access Specifiers in TwinCAT 3}
\end{table}

\begin{figure}[h]
\centering
\begin{tabular}{|c|c|c|c|c|}
\hline
\textbf{Access From} & \textbf{PUBLIC} & \textbf{PRIVATE} & \textbf{PROTECTED} & \textbf{INTERNAL} \\
\hline
Same FB & \checkmark & \checkmark & \checkmark & \checkmark \\
\hline
Child FB & \checkmark & & \checkmark & \checkmark \\
\hline
Same Library & \checkmark & & & \checkmark \\
\hline
External Code & \checkmark & & & \\
\hline
\end{tabular}
\caption{Access Specifier Visibility Matrix}
\end{figure}

% ============================================================
\section{FB Body vs. Methods}
% ============================================================

Unlike traditional languages like Java where all state changes must occur through methods, TwinCAT allows you to call the \textbf{body} of the function block directly.

\begin{itemize}[leftmargin=*]
    \item \textbf{Cyclical Logic (The Body):} Use the FB body for code that must run every single PLC cycle.
    
    \item \textbf{One-Shot Logic (Methods):} Use methods for specific requests or messages to change state that do not need to run every cycle.
\end{itemize}

\begin{table}[h]
\centering
\begin{tabular}{@{}lll@{}}
\toprule
\textbf{Aspect} & \textbf{FB Body} & \textbf{Methods} \\
\midrule
Execution & Every PLC cycle & On-demand (when called) \\
Use Case & Continuous monitoring/control & Specific actions/requests \\
Example & Motor speed regulation loop & \texttt{SetTargetSpeed()} \\
\bottomrule
\end{tabular}
\caption{FB Body vs. Methods}
\end{table}

% ============================================================
\section{Best Practices for Professional Code}
% ============================================================

\subsection{Self-Descriptive Code}

Prioritize descriptive variable and method names over comments. Comments often become outdated as code changes, leading to confusion.

\subsection{Avoid ``Magic Numbers''}

Use \textbf{constants} for fixed values (like buffer sizes) to make the code more readable.

\subsection{Modularization}

Create small function blocks that do ``one thing only'' but do it very well.

\subsection*{Example: Constants vs. Magic Numbers}

\begin{lstlisting}
// BAD PRACTICE: Using "Magic Numbers"
// What does 100 mean? Why 100? Hard to understand and maintain.
FUNCTION_BLOCK FB_EventLogger_Bad
VAR
    astEventBuffer : ARRAY[0..99] OF ST_Event;  // Why 99? What if we need to change it?
    nCurrentIndex  : UDINT;
END_VAR

// Checking against magic number - unclear intent
IF nCurrentIndex < 100 THEN
    // Add event...
END_IF
\end{lstlisting}

\begin{lstlisting}
// GOOD PRACTICE: Using Named Constants
// Clear, self-documenting, easy to maintain
FUNCTION_BLOCK FB_EventLogger_Good
VAR CONSTANT
    MAXIMUM_BUFFER_SIZE : UDINT := 100;    // Single point of definition
    MINIMUM_EVENT_ID    : UDINT := 1;       // Descriptive name explains purpose
    INVALID_INDEX       : UDINT := 16#FFFFFFFF;  // Clear meaning
END_VAR
VAR
    // Array size defined using the constant - change in one place updates everywhere
    astEventBuffer : ARRAY[0..(MAXIMUM_BUFFER_SIZE - 1)] OF ST_Event;
    nCurrentIndex  : UDINT := 0;
    nEventCount    : UDINT := 0;
END_VAR

// Clear intent: checking against named constant
IF nCurrentIndex < MAXIMUM_BUFFER_SIZE THEN
    // Add event...
END_IF
\end{lstlisting}

\begin{warningbox}
Magic numbers make code difficult to maintain. If you need to change a buffer size from 100 to 200, you would have to find and update every occurrence of ``100'' in your code. With constants, you change it in one place.
\end{warningbox}

% ============================================================
\section{Practical Example: The Event Logger}
% ============================================================

The tutorial demonstrates building an \texttt{EventLogger} function block to store a fixed-size buffer of events.

\begin{itemize}[leftmargin=*]
    \item \textbf{Internal Buffer:} An array of event structures.
    
    \item \textbf{Public Method (\texttt{AddEvent}):} Allows external code to submit an event identity, text, etc.
    
    \item \textbf{Private Method (\texttt{IsEventBufferFull}):} An internal helper method used by \texttt{AddEvent} to check if the buffer has reached its maximum capacity before adding new data.
\end{itemize}

\subsection*{Example: Private Helper Method}

\begin{lstlisting}
// Function Block: FB_EventLogger
// A circular buffer for storing system events
FUNCTION_BLOCK FB_EventLogger
VAR CONSTANT
    MAXIMUM_BUFFER_SIZE : UDINT := 100;    // Maximum events to store
END_VAR
VAR
    astEventBuffer : ARRAY[0..(MAXIMUM_BUFFER_SIZE - 1)] OF ST_Event;
    nCurrentIndex  : UDINT := 0;           // Current write position
    nEventCount    : UDINT := 0;           // Total events stored
    bBufferFull    : BOOL := FALSE;        // Buffer status flag
END_VAR
\end{lstlisting}

\begin{lstlisting}
// Private Method: IsEventBufferFull
// Checks if the event buffer has reached maximum capacity
// Access: PRIVATE - can only be called from within FB_EventLogger
METHOD PRIVATE IsEventBufferFull : BOOL
VAR_INPUT
    // No inputs needed - uses internal FB variables
END_VAR

// Compare current index against the maximum buffer size constant
// Returns TRUE if buffer is full, FALSE otherwise
IF nCurrentIndex >= MAXIMUM_BUFFER_SIZE THEN
    IsEventBufferFull := TRUE;
ELSE
    IsEventBufferFull := FALSE;
END_IF

// Alternative one-liner:
// IsEventBufferFull := (nCurrentIndex >= MAXIMUM_BUFFER_SIZE);
\end{lstlisting}

\begin{lstlisting}
// Public Method: AddEvent
// Adds a new event to the buffer if space is available
// Access: PUBLIC - can be called from external code
METHOD PUBLIC AddEvent : BOOL
VAR_INPUT
    nEventId    : UDINT;           // Unique event identifier
    sEventText  : STRING(255);     // Event description
    eSeverity   : TcEventSeverity; // Event severity level
END_VAR
VAR
    fbGetSystemTime : GETSYSTEMTIME;
    stFileTime      : T_FILETIME;
END_VAR

// First check if buffer is full using our private helper method
IF THIS^.IsEventBufferFull() THEN
    // Buffer is full - cannot add new event
    AddEvent := FALSE;
    RETURN;
END_IF

// Buffer has space - add the new event
astEventBuffer[nCurrentIndex].nIdentity      := nEventId;
astEventBuffer[nCurrentIndex].sText          := sEventText;
astEventBuffer[nCurrentIndex].eEventSeverity := eSeverity;

// Get current timestamp
fbGetSystemTime(timeLoDW => stFileTime.dwLowDateTime,
                timeHiDW => stFileTime.dwHighDateTime);
astEventBuffer[nCurrentIndex].dtTimestamp := FILETIME_TO_DT(stFileTime);

// Update index and count
nCurrentIndex := nCurrentIndex + 1;
nEventCount   := nEventCount + 1;

// Return success
AddEvent := TRUE;
\end{lstlisting}

\subsection*{Usage Example}

\begin{lstlisting}
// Program: MAIN
// Demonstrating the use of FB_EventLogger
PROGRAM MAIN
VAR
    fbLogger    : FB_EventLogger;    // Instance of the event logger
    bAddResult  : BOOL;              // Result of AddEvent call
    bTrigger    : BOOL;              // Trigger to add an event
END_VAR

// Add an event when triggered
IF bTrigger THEN
    bAddResult := fbLogger.AddEvent(
        nEventId   := 1001,
        sEventText := 'System started successfully',
        eSeverity  := TcEventSeverity.Info
    );
    bTrigger := FALSE;  // Reset trigger
END_IF
\end{lstlisting}

% ============================================================
\section*{Quick Reference Card}
\addcontentsline{toc}{section}{Quick Reference Card}
% ============================================================

\subsection*{Function Block Declaration Syntax}

\begin{lstlisting}
FUNCTION_BLOCK <Name>
// Or with inheritance:
FUNCTION_BLOCK <Name> EXTENDS <ParentFB>
VAR_INPUT
    // Input parameters
END_VAR
VAR_OUTPUT
    // Output parameters
END_VAR
VAR
    // Internal variables (state)
END_VAR
VAR CONSTANT
    // Named constants
END_VAR

// FB Body code here (runs every cycle when FB is called)
\end{lstlisting}

\subsection*{Method Declaration Syntax}

\begin{lstlisting}
METHOD <AccessSpecifier> <MethodName> : <ReturnType>
VAR_INPUT
    // Input parameters
END_VAR
VAR
    // Local variables
END_VAR

// Method implementation
<MethodName> := <returnValue>;
\end{lstlisting}

\subsection*{OOP Keywords Summary}

\begin{table}[h]
\centering
\begin{tabular}{@{}ll@{}}
\toprule
\textbf{Keyword} & \textbf{Purpose} \\
\midrule
\texttt{EXTENDS} & Inherit from a parent function block \\
\texttt{IMPLEMENTS} & Implement an interface \\
\texttt{THIS\^{}} & Reference to the current instance \\
\texttt{SUPER\^{}} & Reference to the parent class \\
\texttt{METHOD} & Define a method within an FB \\
\texttt{PROPERTY} & Define a property (getter/setter) \\
\texttt{VAR CONSTANT} & Define named constants \\
\bottomrule
\end{tabular}
\caption{OOP Keywords in TwinCAT 3}
\end{table}

\subsection*{Common Naming Conventions}

\begin{table}[h]
\centering
\begin{tabular}{@{}lll@{}}
\toprule
\textbf{Prefix} & \textbf{Type} & \textbf{Example} \\
\midrule
\texttt{FB\_} & Function Block & \texttt{FB\_Motor}, \texttt{FB\_EventLogger} \\
\texttt{I\_} & Interface & \texttt{I\_Notification}, \texttt{I\_Loggable} \\
\texttt{fb} & FB Instance & \texttt{fbWinch}, \texttt{fbLogger} \\
\texttt{UPPER\_CASE} & Constants & \texttt{MAXIMUM\_BUFFER\_SIZE} \\
\bottomrule
\end{tabular}
\caption{Naming Conventions for OOP in TwinCAT 3}
\end{table}

\vfill

\begin{center}
\textit{Last Updated: \today}
\end{center}

\end{document}